\section{The degree test}\label{sec:test}

\subsection{Statement and construction}\label{sec:statement}

We state the guarantee first and give the protocol afterwards. Throughout, $\kappa\ge3$,
$\rho:=k^m/N$ and $\alpha_0:=\sqrt\rho\,(1+\frac1{2\kappa})$; the weights $w\in\N^m$ fix which
linear combination of the degrees the test certifies, and $D$ bounds that combination. The theorem is
stated for an arbitrary hypothesis $\mathrm H(\alpha_*,\Theta)$ of \cref{hyp:ca}:
\cref{prop:H1,prop:H2} supply one at the Johnson radius and in the unique-decoding regime, and
\cref{cor:johnson,cor:ud} give the resulting errors.

\begin{theorem}[Soundness, reduction form]\label{thm:main}
Assume $\mathrm H(\alpha_*,\Theta)$ with $\alpha_*>\rho$, let $\alpha\ge\alpha_*$, and suppose
$\agree(u_0,C_{w,D})\le\alpha$. Then for every prover the test $\Pi_{w,D}$ accepts with
probability at most
\[
\frac1{|\F|}\sum_{i=1}^{m}\Theta\big(|L_i|,k^i,k-1\big)+\alpha^s .
\]
\end{theorem}

The protocol has $m$ folding rounds and $s$ repetitions of the query phase. The bounds on the
format are part of the definition: a word $u_i$ is accepted only if all its values lie in
$\F[T]_{\le d_i}$, and the last message only if $\deg g\le D$.

\begin{construction}[Degree test $\Pi_{w,D}$]\label{def:protocol}
Parameters: $\F,k,m,N,L_0,\dots,L_m$ as in \cref{def:domains}; $w\in\N^m$; $D\in\N$; $s\ge1$.
Put $d_0:=0$ and $d_i:=\min\big(D,(k-1)\sum_{l\le i}w_l\big)$.
\begin{enumerate}
\item The prover sends $u_0:L_0\to\F$ (honestly $u_0=U_f|_{L_0}$).
\item For $i=1,\dots,m$: the verifier sends a uniform $r_i\in\F$; put $c_i:=r_iT^{w_i}$. The
prover sends $u_i:L_i\to\F[T]_{\le d_i}$ (honestly $u_i=\Fold_{c_i}(u_{i-1})$). For $i=m$ the
word $u_m$ is constant, sent as one polynomial $g\in\F[T]_{\le D}$.
\item Repeat $s$ times: pick a uniform $\xi_0\in L_0$, put $\xi_i:=\xi_{i-1}^k$, and check
$u_i(\xi_i)=\Fold_{c_i}(u_{i-1})(\xi_i)$ for $i=1,\dots,m$ (reading $u_{i-1}$ on $\fib(\xi_i)$).
\end{enumerate}
The verifier accepts iff all checks pass.
\end{construction}
% = D8

One repetition of Step~3 reads $u_0$ at one point and then, level by level, the $k$ values of
$u_{i-1}$ on the fibre $\fib(\xi_i)$; consecutive levels share one of these values, because
$\xi_{i-1}\in\fib(\xi_i)$.

\subsection{Completeness}\label{sec:complete}

An honest prover of a polynomial with $\deg_wf\le D$ passes every check for every choice of
challenges, and its last message is the polynomial whose degree is the quantity under test.

\begin{proposition}[Completeness and the terminal polynomial]\label{prop:complete}
Let $f\in\F[x_1,\dots,x_m]_{<k}$ with $\deg_wf\le D$, and let the prover be honest.
\begin{enumerate}
\item[(a)] Every format condition and every query check of $\Pi_{w,D}$ holds for every choice of
challenges; the verifier accepts with probability $1$.
\item[(b)] The final message is $g(T)=f(r_1T^{w_1},\dots,r_mT^{w_m})$.
\item[(c)] $\deg g\le\deg_wf$, and $\Pr_r\big[\deg g<\deg_wf\big]\le\frac{m(k-1)}{|\F|}$.
\end{enumerate}
\end{proposition}
% = P1

\begin{proof}
(a) $u_0=U_f|_{L_0}\in C_{w,D}$. By \cref{lem:stab}(c), each honest $u_i$ takes values of
$T$-degree $\le d_i$, so the formats hold, and each check holds because
$u_i=\Fold_{c_i}(u_{i-1})$ by definition.
(b) is \cref{cor:terminal}.
(c) The coefficient of $T^{\deg_wf}$ in $g$ is $f^{\mathrm{top}}(r)$, where
$f^{\mathrm{top}}:=\sum_{\langle w,\alpha\rangle=\deg_wf}f_\alpha x^\alpha\ne0$. It has total
degree $\le m(k-1)$, so it vanishes with probability $\le m(k-1)/|\F|$ by Schwartz--Zippel, and no
larger power of $T$ occurs.
\end{proof}

\subsection{Query weights}\label{sec:weights}

Following \cite[\S8.2.1]{BCIKS20}, the analysis attaches to each point $y\in L_i$ a weight
$\mu^{(i)}(y)$: the probability that the checks at levels $1,\dots,i$ pass, for a query path
through $y$. The weights are defined level by level from the transcript.

\begin{definition}[Query weights, after {\cite[\S8.2.1]{BCIKS20}}]\label{def:weights}
For a transcript, define $\mu^{(0)}\equiv1$ on $L_0$ and, for $1\le i\le m$ and $y\in L_i$,
\[
\nu^{(i)}(y):=\frac1k\sum_{x\in\fib(y)}\mu^{(i-1)}(x),\qquad
\mu^{(i)}(y):=\begin{cases}\nu^{(i)}(y)&\text{if }u_i(y)=\Fold_{c_i}(u_{i-1})(y),\\0&\text{otherwise.}\end{cases}
\]
The values of $\mu^{(i)},\nu^{(i)}$ are multiples of $1/k^i$, and $k^i|L_i|=N$; hence every
$\mu^{(i)}$- or $\nu^{(i)}$-agreement on $L_i$ is a multiple of $1/N$.
\end{definition}
% = D9

\begin{importedthm}[{\cite[Claim 8.5]{BCIKS20}: acceptance as a weight}]\label{ext:claim85}
With the weights of \cref{def:weights}, the probability that one query of \cref{def:protocol}
accepts equals $\frac1{|L_m|}\sum_{y\in L_m}\mu^{(m)}(y)$.
\end{importedthm}
% = E2

The proof in \cite{BCIKS20} samples $y\in L_m$ uniformly and then a uniform element of each fibre
downwards; sampling a uniform $\xi_0\in L_0$ and setting $\xi_i:=\xi_{i-1}^k$ gives the same
distribution of query paths, and their induction on $i$ (p.~43) applies verbatim.

\subsection{The round lemma}\label{sec:round}

Fix a round and a word that is far from the graded code. The lemma bounds the number of
challenges for which the folded word becomes close, and it is the only place where the hypothesis
of \cref{sec:params} enters the estimate.

\begin{lemma}[Round lemma]\label{lem:round}
Assume $\mathrm H(\alpha_*,\Theta)$ and $\alpha_*>\rho$. Let $1\le i\le m$, let $u:L_{i-1}\to\F[T]_{\le d_{i-1}}$,
and let $\mu^{(i-1)}$ be a weight on $L_{i-1}$ with values in $\frac1{k^{i-1}}\Z$; let $\nu^{(i)}$
be its fibre average (\cref{def:weights}). Put
\[
\beta_i:=\max\Big(\agree_{\mu^{(i-1)}}\big(u,\Gc_{i-1}^{w,D}\big),\ \alpha_*\Big),\qquad
S:=\Big\{r\in\F:\ \agree_{\nu^{(i)}}\big(\Fold_{rT^{w_i}}(u),\Gc_i^{w,D}\big)>\beta_i\Big\}.
\]
Then $|S|\le\Theta\big(|L_i|,k^i,k-1\big)$.
\end{lemma}
% = L13

\begin{proof}
\begin{enumerate}
\item Assume $|S|>\Theta(|L_i|,k^i,k-1)$, let $u_j$ be the components of $u$ and put
$\tilde u_j:=T^{w_ij}u_j$; by \cref{def:fold}, $\Fold_{rT^{w_i}}(u)=\sum_{j<k}r^j\tilde u_j$ for
every $r\in\F$.
\item By \cref{lem:gran}, the set $S$ equals
$\{r:\agree_{\nu^{(i)}}(\sum_jr^j\tilde u_j,\Gc_i^{w,D})\ge\beta_i^+\}$, since $\nu^{(i)}$ has
denominator $k^i$, and $\beta_i^+\ge\alpha_*$.
\item By \cref{lem:ext} with $\E\supseteq\F$ of degree $t:=1+\max(D,d_{i-1}+w_i(k-1))$, the
$\tilde u_j$ become words over $\E$ and $\Gc_i^{w,D}$ an $\F$-linear subspace $W$ of
$\RS[\E,L_i,k^{m-i}]$, of rate $\rho$.
\item So \cref{cor:share} with $n=|L_i|$, $M=k^i$ and $l=k-1$ gives $v_0,\dots,v_{k-1}\in\Gc_i^{w,D}$
and $A\subseteq L_i$ with $\nu^{(i)}(A)\ge\beta_i^+$ and $\tilde u_j=v_j$ on $A$.
\item $|A|\ge\nu^{(i)}(A)|L_i|>\rho|L_i|=k^{m-i}$ and $T^{w_ij}\mid v_j$ on $A$, so \cref{lem:div}
gives $v_j=T^{w_ij}p_j$ with $p_j\in\Gc_i^{w,D-w_ij}$; then $u_j=p_j$ on $A$.
\item By \cref{lem:unfold}, $h:=\sum_{j<k}x^jp_j(x^k)$ lies in $\Gc_{i-1}^{w,D}$ with components
$h_j=p_j$.
\item By \cref{lem:fold-sub}(c), $u(x)=\sum_jx^ju_j(x^k)=h(x)$ for every $x\in\fib(y)$, $y\in A$.
\item By \cref{lem:lift}, $\agree_{\mu^{(i-1)}}(u,h)\ge\nu^{(i)}(A)\ge\beta_i^+>\beta_i$,
contradicting the definition of $\beta_i$.
\end{enumerate}
\end{proof}

\subsection{Proof of \texorpdfstring{Theorem~\ref{thm:main}}{the main theorem}}\label{sec:proof-main}

The proof collects the rounds in which the challenge is bad and then counts the probability that a
query survives all of them.

\begin{proof}[Proof of Theorem~\ref{thm:main}]
\begin{enumerate}
\item By \cref{lem:start}, $\agree(u_0,\Gc_0^{w,D})=\agree(u_0,C_{w,D})\le\alpha$.
\item For $1\le i\le m$ let $B_i$ be the event that $r_i$ lies in the set $S_i$ of
\cref{lem:round}; then $\Pr[B_i]\le\Theta(|L_i|,k^i,k-1)/|\F|$, so $\Pr[\bigcup_iB_i]$ is at most
the first term.
\item By \cref{lem:no-increase}, outside $B_i$ one has
$\agree_{\mu^{(i)}}(u_i,\Gc_i^{w,D})\le\max\big(\agree_{\mu^{(i-1)}}(u_{i-1},\Gc_{i-1}^{w,D}),\alpha_*\big)$,
so induction from step~1 gives $\agree_{\mu^{(m)}}(u_m,\Gc_m^{w,D})\le\alpha$.
\item Since $u_m=g\in\Gc_m^{w,D}$ and $\Gc_m^{w,D}$ is the set of constant words with value in
$\F[T]_{\le D}$, \cref{ext:claim85} bounds the probability that one query accepts by
$\agree_{\mu^{(m)}}(g,g)\le\alpha$, and the $s$ repetitions by $\alpha^s$.
\end{enumerate}
\end{proof}

\begin{remark}[Characteristic]\label{rem:char}
\Cref{prop:H1} inherits from \cref{ext:bciks72} the separability convention of
\cite[\S5.2.3]{BCIKS20}: the characteristic of $\F$ exceeds the bound $D_Y$ of
\cref{ext:proof-facts}(F2);
the general case, with inseparable factors, is treated in \cite[Appendix~C]{BCIKS20}.
\end{remark}

\begin{corollary}[Soundness at the Johnson radius]\label{cor:johnson}
Let $\kappa\ge3$ be an integer, $\rho:=k^m/N$, $\alpha\ge\alpha_0:=\sqrt\rho\,(1+\frac1{2\kappa})$, and
suppose $\agree(u_0,C_{w,D})\le\alpha$. Then the test accepts with probability at most
$\varepsilon_C+\alpha^s$, where
\[
\varepsilon_C:=\frac1{|\F|}\Big(\frac{(1+\frac1{2\kappa})^7\kappa^7\,N^2}{3(k+1)\rho^{3/2}}
+m(k-1)\frac{(2\kappa+1)(N+1)}{\sqrt\rho}\Big).
\]
\end{corollary}
% = C2

\begin{proof}
Use \cref{prop:H1} in \cref{thm:main}: for each $i$,
$\Theta_J(|L_i|,k^i,k-1)=(k-1)\max(A|L_i|^2,B(N+1))$, with $A$ and $B$ the two constants of
$\Theta_J$ and $M=k^i$, so $Mn=N$. Hence
\[
\sum_{i=1}^{m}\Theta_J(|L_i|,k^i,k-1)\ \le\ A\sum_{i=1}^{m}(k-1)|L_i|^2+B\,m(k-1)(N+1),
\]
and $\sum_i(k-1)|L_i|^2=(k-1)N^2\sum_ik^{-2i}=\frac{N^2}{k+1}$.
\end{proof}

\begin{corollary}[Soundness in the unique-decoding regime]\label{cor:ud}
Let $\alpha>\frac{1+\rho}2$ and suppose $\agree(u_0,C_{w,D})\le\alpha$. Then the test accepts with
probability at most $\varepsilon_{\mathrm U}+\alpha^s$, where
\[
\varepsilon_{\mathrm U}:=\frac{m(k-1)(N+1)}{|\F|}\Big(\frac1{\eta_{\mathrm U}}+\frac3{\sqrt\rho}\Big),
\qquad \eta_{\mathrm U}=\min\Big(\frac{(1-\sqrt\rho)^2}2,\frac{\sqrt\rho}{20}\Big).
\]
\end{corollary}
% = C3

\begin{proof}
Use \cref{prop:H2} with $\alpha_*=\alpha$ in \cref{thm:main}: each of the $m$ terms
$\Theta_U(|L_i|,k^i,k-1)$ equals $(k-1)(N+1)(\frac1{\eta_{\mathrm U}}+\frac3{\sqrt\rho})$, since
$k^i|L_i|=N$.
\end{proof}

\subsection{Parameters}\label{sec:numbers}

We instantiate Corollary~\ref{cor:johnson} at a concrete point. Take $N=2^{20}$, $k=4$, $m=8$,
$w=(1,\dots,1)$ and $D=6$, so $\rho=k^m/N=2^{16-20}=2^{-4}$, together with $|\F|=2^{256}$, which
satisfies $N\mid|\F|-1$, and $\kappa=2^{11}-1$. Corollary~\ref{cor:johnson} bounds the error by
$\varepsilon_C+\alpha^s$ with $\alpha\ge\alpha_0$, and this bound is smallest at
$\alpha=\alpha_0$.

\begin{center}
\renewcommand{\arraystretch}{1.15}
\begin{tabular}{@{}p{6.4cm}lr@{}}
\hline
quantity & value & $\log_2$ \\
\hline
$\kappa+\tfrac12$ & $2047.5$ & $11.0$ \\
$(\kappa+\tfrac12)^7$ & $1.509\cdot10^{23}$ & $77.0$ \\
$N^2/(3(k+1)\rho^{3/2})$ & $4.691\cdot10^{12}$ & $42.1$ \\
first summand of $\varepsilon_C$ & $7.077\cdot10^{35}$ & $119.1$ \\
second summand of $\varepsilon_C$ & $4.122\cdot10^{11}$ & $38.6$ \\
$\varepsilon_C$ & $6.112\cdot10^{-42}$ & $-136.9$ \\
$\alpha_0$ & $0.2500611$ & $-2.0$ \\
$\alpha_0^{65}$ & $1.047\cdot10^{-39}$ & $-130.0$ \\
\hline
\end{tabular}
\end{center}

The second summand is $m(k-1)(2\kappa+1)(N+1)/\sqrt\rho=24\cdot4095\cdot4\cdot(2^{20}+1)$,
and the first is $(1+\frac1{2\kappa})^7\kappa^7=(\kappa+\frac12)^7$ times
$N^2/(3(k+1)\rho^{3/2})=2^{46}/15$. With $s=65$ repetitions,
\[
\varepsilon_C+\alpha^s\le 2^{-136.9}+2^{-130.0}<2^{-129.9}.
\]
The query term dominates: the exponent is set by $s\log_2\alpha_0$, and each additional
repetition lowers it by $2.0$.

At the same instance Corollary~\ref{cor:ud} covers the smaller radius. Here
$\eta_{\mathrm U}=\min\{0.125,\,0.0125\}=0.0125$, so $1/\eta_{\mathrm U}+3/\sqrt\rho=92$ and
\[
\varepsilon_{\mathrm U}=\frac{24\cdot(2^{20}+1)\cdot92}{2^{256}}=2^{-224.9}
\]
for every $\alpha>\frac{1+\rho}2=0.625$, that is, up to distance $0.375$.

A repetition of the query phase at this instance reads $k\sum_{i=1}^{m}(d_{i-1}+1)=4\,(1+4+6\cdot7)=188$ field
elements per repetition, below the bound $km(D+1)=224$; FRI on the same word reads
$km=32$.
